% 第 2 章 RedHawk Flow（原书 2-5 … 2-8）
\bichapter{RedHawk 流程}{RedHawk Flow}
\label{chap:flow}

\section{引言}
\subsection*{Introduction}

RedHawk 可对电路执行多种电源分析：
\begin{itemize}
  \item 基于平均周期电流的静态（IR）压降分析
  \item 基于最坏开关电流的动态压降分析
  \item 电迁移（electromigration）分析
  \item 关键路径与时钟树影响分析
\end{itemize}

每类分析可按可用输入数据、期望的分析速度与结果精度，以不同方式运行。以下各节概述静态 IR 与动态压降分析的数据流。

\section{静态压降分析流程}
\subsection*{Static Voltage Drop Analysis Flow}

图~\ref{fig:rh-s-flow} 给出运行 RedHawk-S（静态 IR drop 方案）的设计流程。静态压降分析的关键步骤如下：

\begin{enumerate}
  \item \textbf{准备设计数据与输入文件}（见第~\ref{chap:ui}~章）。
  \begin{itemize}
    \item 准备描述 IC 工艺的 RedHawk technology 文件数据（如 \texttt{tutorial.tech}）。
    \item 准备 pad 单元名、pad 实例名或 pad 位置文件（\texttt{tutorial.pcell}、\texttt{tutorial.pad} 和/或 \texttt{tutorial.ploc}）。
    \item 生成包含 slew、timing window 与时钟实例的 STA 输出文件（\texttt{tutorial.sta}）。参见第~\ref{chap:ate}~章「使用 Apache Timing Engine（ATE）创建时序文件」。
    \item 准备 Global System Requirements（GSR）文件（引用 \texttt{.tech} 文件、pad 文件、STA 文件、LEF 文件、DEF 文件、LIB 文件等），用于静态 IR/EM 和/或动态压降分析（\texttt{tutorial.gsr}）。
    \item 使用 GSR 文件导入设计数据（\texttt{tutorial.gsr}）。
  \end{itemize}
  \item 由 \texttt{.lib} 单元数据执行功耗计算，或导入先前已计算的功耗数据（见第~\ref{chap:static}~章）。
  \item 提取电源网格（R 网络）。
  \item 执行静态 IR/EM 分析（见第~\ref{chap:static}~章）。
  \item 生成并查看 IR/EM 结果的地图与文本报告（见第~\ref{chap:reports}~章）。
  \item 进行 what-if 分析，并通过网格修改与优化修复关键静态 IR drop 区域（见第~\ref{chap:grid}~章）。
\end{enumerate}

\figplaceholder{Figure 2-1}{RedHawk-S 静态 IR 功耗分析流程}
\label{fig:rh-s-flow}

图~\ref{fig:rh-s-flow} 所示流程的输入包括：工艺相关库（LEF/LIB 文件、Apache Tech 文件）、物理设计（DEF/GDS 文件）、设计与分析规格（GSR 文件、封装规格、\texttt{.pcell}/\texttt{.ploc}/\texttt{.pad} 文件），以及可选的更高精度寄生负载（SPEF/DSPF 文件）与实例 slew（STA 文件）。

RedHawk-S 输出包括：
\begin{itemize}
  \item IR 压降等高线/色阶图（contour maps）
  \item 电迁移（EM）分析
  \item 功耗密度与平均电流图
  \item 详细静态功耗、电压与电流数据的文本报告文件
  \item 警告与违例报告
\end{itemize}

\section{动态压降分析流程}
\subsection*{Dynamic Voltage Drop Analysis Flow}

图~\ref{fig:rh-ev-flow} 给出运行 RedHawk-EV（动态压降方案）的设计流程。

\figplaceholder{Figure 2-2}{RedHawk-EV 动态功耗分析流程}
\label{fig:rh-ev-flow}

动态分析流程的关键步骤如下：

\begin{enumerate}
  \item \textbf{准备设计数据与输入文件}（见第~\ref{chap:ui}~章）。
  \begin{noteBox}
  若已运行过 RedHawk 静态 IR drop 分析，则除在 GSR 中设置特定动态运行参数外，本步无需重复。
  \end{noteBox}
  \begin{itemize}
    \item 准备描述 IC 工艺的 RedHawk technology 文件数据（如 \texttt{tutorial.tech}）。
    \item 准备 pad 单元名、pad 实例名或 pad 位置文件（\texttt{tutorial.pcell}、\texttt{tutorial.pad} 和/或 \texttt{tutorial.ploc}）。
    \item 生成包含 slew、timing window 与时钟实例的 STA 输出文件。参见第~\ref{chap:ate}~章「使用 Apache Timing Engine（ATE）创建时序文件」。
    \item 准备 Global System Requirements（GSR）文件（引用 \texttt{.tech} 文件、pad 文件、STA 文件、LEF 文件、DEF 文件、LIB 文件等），用于动态压降分析（\texttt{tutorial.gsr}）（文件定义见附录~\ref{chap:files}「文件定义」）。
    \item 使用 GSR 文件导入设计数据（\texttt{tutorial.gsr}）。
  \end{itemize}
  \item \textbf{准备 RedHawk-EV 额外所需输入}（相对静态分析，见第~\ref{chap:ui}~章「用户界面与数据准备」）：
  \begin{itemize}
    \item 来自 STA 的 timing window 与 slew（推荐）
    \item 来自 SPEF 或 DSPF 的提取寄生参数（推荐）
    \item Pad、wirebond/bump 与封装 R、L、C、K 信息
    \item 工艺数据——导体厚度、介质厚度与介电常数（见附录~\ref{chap:files}「Apache Technology File (\texttt{*.tech})」）
    \item SPICE model card 与库 subcircuit（表征 APL 电流波形所必需）
    \item 全部存储器、I/O 与 IP block 的 SPICE subcircuit（可选）
  \end{itemize}
  \item 由 \texttt{.lib} 单元数据计算详细功耗分布，或导入先前已计算的功耗数据（见第~\ref{chap:static}~章）。
  \item 提取电源网格（RLC 网络）。
  \item 运行 APL（Apache Power Library）表征，获得典型 corner 条件下的电流剖面（current profiles）、电源回路的有效串联电阻（ESR），以及去耦电容与漏电流（见第~\ref{chap:apl}~章「使用 Apache Power Library 进行表征」）。
  \item 执行动态压降与峰值电流分析（见第~\ref{chap:dvd}~章「动态压降分析」）。
  \item 生成并查看动态分析结果的地图与文本报告（见第~\ref{chap:reports}~章「报告」）。
  \item 进行 what-if 分析，优化 decap 布局、修改网格并执行实例替换，以修复热点——动态压降最高的区域（见第~\ref{chap:grid}~章「网格与电源性能修复与优化」）。
\end{enumerate}

图~\ref{fig:rh-ev-flow} 所示流程在静态分析输入基础上，额外需要 Apache Power Library 中的电流剖面、实例功耗电阻与 intrinsic decap；RedHawk-EV 执行功耗计算、电源网格 RLC 提取、动态压降热点分析与修复，以及 Vectorless Dynamic 仿真、封装寄生参数与时序影响分析。

RedHawk-EV 输出包括：
\begin{itemize}
  \item 动态压降色阶图与功耗密度色阶图（见第~\ref{chap:dvd}~章「动态压降分析」）
  \item 电容图（含 decap 效应，见第~\ref{chap:dvd}~章「动态压降分析」）
  \item 报告文件（见第~\ref{chap:reports}~章「报告」）
  \item 面向布局布线环境的 ECO 脚本（见第~\ref{chap:grid}~章「网格与电源性能修复与优化」）
\end{itemize}
